\subsection{Elementary homological properties of regular local rings}

PROPOSITION 1.- Let A be a regular Noetherian local ring and n its dimension. Then we have \(\mathrm{dh(A)} = n\) and, for every integer \(i \geqslant 0\) ,

\[
\left[ \operatorname{Ext} _ {A} ^ {i} \left(\kappa_ {A}, \kappa_ {A}\right): \kappa_ {A} \right] = \left[ \operatorname{Tor} _ {i} ^ {A} \left(\kappa_ {A}, \kappa_ {A}\right): \kappa_ {A} \right] = \binom {n} {i}.
\]

Let \(\mathbf{x} = (x_{1},\ldots ,x_{n})\) a coordinate system of A (VIII, § 5, no. 1, def. 1). The sequence \(\mathbf{x}\) generate \(\mathfrak{m}_{\mathrm{A}}\) and is completely secant for A (loc. cit., no. 2, th. 1). The Koszul complex \(\mathbf{K}_{\bullet}(\mathbf{x},\mathrm{A})\) is a free resolution of \(\kappa_{\mathrm{A}}\) (A, X, p.~159, remark 3), whose differential is zero modulo \(\mathfrak{m}_{\mathrm{A}}\) . For every integer \(i\geqslant 0\) , we therefore have ( \(\S 3\) , no. 3, formula (1))

\[
\left[ \operatorname{Ext} _ {A} ^ {i} (\kappa_ {A}, \kappa_ {A}): \kappa_ {A} \right] = \left[ \operatorname{Tor} _ {i} ^ {A} (\kappa_ {A}, \kappa_ {A}): \kappa_ {A} \right] = \operatorname{rg} _ {A} (\mathbf {K} _ {i} (\mathbf {x}, A)) = \binom {n} {i}.
\]

It then follows from cor. 1 of prop. 4 of no. 3, § 3 that dh(A) = n.

PROPOSITION 2.--- A regular noetherian local ring is factorial.

By prop. 1, every finitely generated module over a regular noetherian local ring admits a projective resolution of finite length by finitely generated projective modules, hence free modules (II, § 3, no. 2, cor. 2 of prop. 5). It then follows from VII, § 4, no. 7, cor. 3 of prop. 16 that such a ring is factorial.

PROPOSITION 3.- Let A be a regular noetherian local ring and M a nonzero finitely generated A-module. Its projective dimension is finite and one has

\[
\mathrm{dp} _ {\mathrm{A}} (\mathrm{M}) + \operatorname{prof} _ {\mathrm{A}} (\mathrm{M}) = \dim (\mathrm{A}).
\]

Indeed, M has finite projective dimension (prop. 1), and one has \(\operatorname{prof}(A) = \dim(A)\) since A is a Macaulay ring ( \(\S\) 2, no. 5, example 7). One then applies th. 1 of \(\S\) 3, no. 5.

COROLLARY 1. One has \(\mathrm{dp}_{\mathrm{A}}(\mathrm{M}) \geqslant \dim(\mathrm{A}) - \dim(\mathrm{M})\) ; for equality to hold, it is necessary and sufficient that M be Macaulay.

COROLLARY 2.--- For the A-module M to be free, it is necessary and sufficient that it be Macaulay and of dimension \(\dim(A)\) or again that it be of depth \(\geqslant \dim(A)\) .

COROLLARY 3.--- Every finitely generated reflexive module over a regular noetherian local ring of dimension 2 is free.

Indeed, a regular noetherian local ring is integrally closed (VIII, § 5, no. 2, cor. 1 of th. 1). Corollary 3 therefore follows from Corollary 2 and § 1, no. 10, prop. 16.

COROLLARY 4.--- Let \(\rho: A \to B\) a local homomorphism of noetherian local rings. One assumes that A is regular and that \(\rho\) makes B a finitely generated A-module. One then has \(\mathrm{dp}_{\mathrm{A}}(\mathrm{B}) \geqslant \dim(\mathrm{A}) - \dim(\mathrm{B})\) . For B to be a Macaulay ring, it is necessary and sufficient that one have \(\mathrm{dp}_{\mathrm{A}}(\mathrm{B}) = \dim(\mathrm{A}) - \dim(\mathrm{B})\) . For B to be a Macaulay ring of dimension equal to \(\dim(\mathrm{A})\) , it is necessary and sufficient that the A-module B be free.

Indeed, one has \(\dim(B) = \dim_A(B)\) (VIII, § 2, no. 3, th. 1); moreover, B is a Macaulay ring if and only if it is a Macaulay A-module (§ 2, no. 6, prop. 8). It therefore suffices to apply Corollaries 1 and 2.

Remark.--- Corollary 4 allows one to characterize Macaulay local rings in several important cases. Let A be a noetherian local ring; it is a Macaulay ring if and only if its completion \(\widehat{\mathsf{A}}\) is (§ 2, no. 7, cor. 4 of prop. 9). Suppose henceforth that the local ring A is complete and set \(d = \dim(\mathsf{A})\) .

\begin{enumerate}
\def\labelenumi{\alph{enumi})}
\tightlist
\item
  Suppose that A possesses a subfield. It then possesses a coefficient subfield K (IX, § 3, no. 3), and there exists a formal power series algebra E = K{[}{[}T₁, \ldots, Tₙ{]}{]} and a surjective homomorphism of K-algebras E → A (loc. cit.); there also exists a formal power series algebra E' = K{[}{[}T₁, \ldots, T\_d{]}{]} and an injective local homomorphism of K-algebras E' → A such that A is a finite algebra over E' (loc. cit.). The following properties are equivalent:
  \begin{enumerate}
  \def\labelenumii{(\roman{enumii})}
  \item
    A is a Macaulay ring;
  \item
    ; we have \(\mathrm{dp}_{\mathrm{E}}(\mathrm{A}) = n - d\)
  \item
    A is a free E'-module.
  \end{enumerate}

\item
  Suppose that the residue field of A has characteristic p \textgreater{} 0. There exists a p-ring of length +∞, with residue field \(\kappa_{\Lambda}\) (IX, § 2, no. 3, prop. 5). Let C be such a ring; there exists a formal power series algebra \(E = C[[T_{1}, \ldots, T_{n}]]\) and a surjective homomorphism \(\rho : E \to \Lambda\) (IX, § 2, no. 5, th. 3). The following properties are equivalent:
  \begin{enumerate}
  \def\labelenumii{(\roman{enumii})}
  \item
    A is a Macaulay ring;
  \item
    ; we have \(\mathrm{dp}_{\mathbf{E}}(\mathrm{A}) = n + 1 - d\)
  \end{enumerate}

Suppose moreover that \(p1_{\Lambda}\) is not a zero divisor in A; there then exists a formal power series algebra \(E' = K[[T_1, \ldots, T_{d-1}]]\) and an injective local homomorphism of K-algebras \(E' \to A\) such that A is a finite algebra over \(E'\) (loc. cit.). The local ring \(E'\) is regular, of dimension \(n+1\) (VIII, § 5, n° 5, example 2). The preceding conditions are also equivalent to

  \begin{enumerate}
  \def\labelenumii{(\roman{enumii})}
  \setcounter{enumii}{2}
  \tightlist
  \item
    A is a free E'-module.
  \end{enumerate}
\end{enumerate}

For analogous results in the case of modules, see § 5, n° 5.

